\section{Conclusion}

Nous avons, dans le présent document, exploré la théorie des transformées de Fourier, formalisé la notion de transformée de Fourier discrète et mis en \oe uvre différents algorithmes qui en découlent dans le domaine du traitement d'image.

Nous avons ainsi pu comprendre le concept de spectre d'une image et comment cette représentation des données pouvait être exploitée dans le cadre de la compression d'image. Nous avons alors pu constater le gain important d'efficacité permit par les algorithmes FFT, passant de $\Theta \left( n^2 \right)$ à $\Theta \left( n \log n \right)$ ainsi que les optimisations successives de l'algorithme de Cooley-Tukey.

Ces méthodes particulièrement optimisées ont alors été employées pour la compression d'image et nous avons une fois encore pu mettre en place différents algorithmes. Nous avons dans le même temps étudiés plusieurs métriques permettant d'évaluer ces algorithmes.

Cette comparaison nous a permis de mettre en lumière la supériorité, en terme de qualité d'image à ratio de compression équivalent, de l'approche du format JPEG utilisant la DCT alliée à un découpage par blocs par rapport à une simple FFT.